%---------------------------------------------------------------------
% UNIT 6 --- DETERMINANTS AND INVERSES OF MATRICES
% Source: Section 6.8 Self Assessment Questions, pp. 131-132.
%---------------------------------------------------------------------
\chapter{Determinants and Inverses}

\srcnote{Source: \emph{Business Mathematics} 5405/1429, \S6.8 \saq{Self Assessment Questions}, pp.~131--132.}

\begin{keyidea}[Cofactors, inverse, Cramer]
For a $3\times3$ matrix, the cofactor of entry $a_{ij}$ is
$C_{ij}=(-1)^{i+j}M_{ij}$, where $M_{ij}$ is the determinant left after
deleting row $i$ and column $j$. The sign pattern is
\[
\mat{+&-&+\\-&+&-\\+&-&+}.
\]
Then
\[
\det A=\sum_j a_{ij}C_{ij}\ \text{(any row)},\qquad
A^{-1}=\frac{1}{\det A}\,\operatorname{adj}A,\qquad
\operatorname{adj}A=\bigl[C_{ij}\bigr]^{t}.
\]
\textbf{Cramer's rule:} for $Ax=b$, $\;x_k=\dfrac{D_k}{D}$, where $D=\det A$ and
$D_k$ replaces column $k$ of $A$ by $b$. It requires $D\neq0$.
\end{keyidea}

\begin{keyidea}[Four shortcuts that turn a long determinant into one line]
\begin{itemize}
  \item Expand along the row or column with the \emph{most zeros}.
  \item A common factor in a whole row or column comes outside the determinant.
  \item If two rows (or columns) are \emph{identical} or \emph{proportional},
        $\det=0$.
  \item Adding a multiple of one row to another leaves the determinant
        unchanged --- use it to manufacture zeros.
\end{itemize}
\end{keyidea}

%---------------------------------------------------------------------
\section{Evaluating determinants}

\begin{question}{Q.1}
Evaluate the following determinants, using properties of determinants where
possible.
\[
\text{(a) }\dt{2&3&-1\\1&1&0\\2&-3&5}\quad
\text{(b) }\dt{2a&a&a\\b&2b&b\\c&c&2c}\quad
\text{(c) }\dt{1&2&3x\\2&3&6x\\3&5&9x}
\]
\[
\text{(d) }\dt{a-b&b-c&c-a\\b-c&c-a&a-b\\c-a&a-b&b-c}\quad
\text{(e) }\dt{2&3&-1&1\\3&0&2&5\\0&8&6&3\\7&-4&5&0}\quad
\text{(f) }\dt{1&0&-2&3\\0&0&0&-5\\4&2&1&4\\1&5&-4&3}
\]
\end{question}

\begin{solution}
\textbf{(a)} Expand along the second row, which contains a zero:
\[
\begin{aligned}
\dt{2&3&-1\\1&1&0\\2&-3&5}
&=2\dt{1&0\\-3&5}-3\dt{1&0\\2&5}+(-1)\dt{1&1\\2&-3}\\
&=2(5-0)-3(5-0)-1(-3-2)\\
&=10-15+5=0 .
\end{aligned}
\]
A zero determinant means the rows are linearly dependent --- indeed
$R_3=3R_2-R_1$.

\medskip
\textbf{(b)} Factor $a$ from row 1, $b$ from row 2, $c$ from row 3:
\[
\dt{2a&a&a\\b&2b&b\\c&c&2c}=abc\dt{2&1&1\\1&2&1\\1&1&2}.
\]
\[
\dt{2&1&1\\1&2&1\\1&1&2}=2(4-1)-1(2-1)+1(1-2)=6-1-1=4 .
\]
\[
\text{Value}=4abc .
\]

\medskip
\textbf{(c)} The third column is $3x$ times $(1,2,3)^{t}$. Factor $3x$ out:
\[
\dt{1&2&3x\\2&3&6x\\3&5&9x}=3x\dt{1&2&1\\2&3&2\\3&5&3}.
\]
Columns 1 and 3 are now identical, so the determinant is 0:
\[
\text{Value}=3x\cdot 0=0 .
\]

\medskip
\textbf{(d)} Each row sums to zero:
$(a-b)+(b-c)+(c-a)=0$, and likewise for rows 2 and 3. So adding columns 2 and 3
to column 1 --- an operation that does not change the determinant --- makes
column 1 entirely zero:
\[
C_1\to C_1+C_2+C_3:\qquad
\dt{0&b-c&c-a\\0&c-a&a-b\\0&a-b&b-c}=0 .
\]
\[
\text{Value}=0 .
\]

\medskip
\textbf{(e)} A $4\times4$. Expand along row 3, $(0,\,8,\,6,\,3)$, whose leading
zero kills one term. With signs $(-1)^{3+j}$ giving $-,+,-,+$ across the row,
so that positions 2 and 4 carry a minus:
\[
\det = -8M_{32}+6M_{33}-3M_{34},
\]
\[
M_{32}=\dt{2&-1&1\\3&2&5\\7&5&0}
=2(0-25)+1(0-35)+1(15-14)=-50-35+1=-84,
\]
\[
M_{33}=\dt{2&3&1\\3&0&5\\7&-4&0}
=2(0+20)-3(0-35)+1(-12-0)=40+105-12=133,
\]
\[
M_{34}=\dt{2&3&-1\\3&0&2\\7&-4&5}
=2(0+8)-3(15-14)-1(-12-0)=16-3+12=25 .
\]
\[
\det=-8(-84)+6(133)-3(25)=672+798-75=1395 .
\]
\vcheck{Expanding instead along column 1 gives
$2(296)+3(-31)+7(128)=592-93+896=1395$. Two independent expansions agree.}

\medskip
\textbf{(f)} Row 2 is $(0,0,0,-5)$ --- three zeros, so expand there and only one
term survives. The sign at position $(2,4)$ is $(-1)^{2+4}=+$:
\[
\det=(-5)\,M_{24}, \qquad
M_{24}=\dt{1&0&-2\\4&2&1\\1&5&-4}.
\]
\[
M_{24}=1(2(-4)-1(5))-0+(-2)(4(5)-2(1))=1(-13)-2(18)=-13-36=-49 .
\]
\[
\det=(-5)(-49)=245 .
\]

\ans{(a) $0$ \quad (b) $4abc$ \quad (c) $0$ \quad (d) $0$ \quad
(e) $1395$ \quad (f) $245$}
\end{solution}

\begin{pitfall}
Parts (a), (c) and (d) are all zero, and each for a different structural reason
--- dependent rows, a repeated column after factoring, rows summing to zero.
Recognising the structure takes one line; brute-force expansion takes ten and
usually goes wrong. Look for the pattern before you expand.
\end{pitfall}

%---------------------------------------------------------------------
\section{Inverses, and whether they distribute over a product}

\begin{question}{Q.2}
Let $A=\mat{1&3&-1\\5&2&3\\-2&1&0}$ and $B=\mat{-1&-2&0\\3&5&1\\-2&-3&4}$.
\textbf{(a)} Find $A^{-1}$, $B^{-1}$, $A^{-1}B^{-1}$, $B^{-1}A^{-1}$ and
$(AB)^{-1}$.
\textbf{(b)} Is it true that $(AB)^{-1}=A^{-1}B^{-1}$?
\textbf{(c)} Is it true that $(AB)^{-1}=B^{-1}A^{-1}$?
\textbf{(d)} Does $A+B$ have an inverse? If so, what is it?
\textbf{(e)} Does $A-B$ have an inverse? If so, what is it?
\end{question}

\begin{solution}
\textbf{(a) $A^{-1}$.} First the determinant, expanding along row 1:
\[
\det A=1(0-3)-3(0+6)+(-1)(5+4)=-3-18-9=-30 \neq 0 .
\]
Cofactors of $A$, with the $+,-,+$ sign pattern applied:
\[
\begin{array}{lll}
C_{11}=-3, & C_{12}=-6, & C_{13}=9,\\
C_{21}=-1, & C_{22}=-2, & C_{23}=-7,\\
C_{31}=11, & C_{32}=-8, & C_{33}=-13 .
\end{array}
\]
The adjoint is the \emph{transpose} of the cofactor matrix:
\[
\operatorname{adj}A=\mat{-3&-1&11\\-6&-2&-8\\9&-7&-13},
\qquad
A^{-1}=\frac{1}{-30}\mat{-3&-1&11\\-6&-2&-8\\9&-7&-13}
=\frac{1}{30}\mat{3&1&-11\\6&2&8\\-9&7&13}.
\]
\vcheck{Row 1 of $A$ times column 1 of $A^{-1}$:
$\tfrac{1}{30}\bigl(1(3)+3(6)+(-1)(-9)\bigr)=\tfrac{30}{30}=1$. Row 1 times
column 3: $\tfrac{1}{30}(-11+24-13)=0$. Correct.}

\medskip
\textbf{$B^{-1}$.}
\[
\det B=-1(20+3)+2(12+2)+0=-23+28=5\neq 0 .
\]
\[
\operatorname{adj}B=\mat{23&8&-2\\-14&-4&1\\1&1&1},
\qquad
B^{-1}=\frac{1}{5}\mat{23&8&-2\\-14&-4&1\\1&1&1}.
\]
\vcheck{Row 1 of $B$ times column 1 of $B^{-1}$:
$\tfrac15\bigl(-1(23)+(-2)(-14)+0(1)\bigr)=\tfrac{5}{5}=1$.}

\medskip
\textbf{$AB$ and its inverse.}
\[
AB=\mat{10&16&-1\\-5&-9&14\\5&9&1},\qquad
\det(AB)=\det A\cdot\det B=(-30)(5)=-150 .
\]
\[
B^{-1}A^{-1}=\frac{1}{5}\mat{23&8&-2\\-14&-4&1\\1&1&1}\cdot
\frac{1}{-30}\mat{-3&-1&11\\-6&-2&-8\\9&-7&-13}
=\frac{1}{-150}\mat{-135&-25&215\\75&15&-135\\0&-10&-10},
\]
\[
B^{-1}A^{-1}=\mat{\tfrac{9}{10}&\tfrac16&-\tfrac{43}{30}\\[3pt]
-\tfrac12&-\tfrac1{10}&\tfrac{9}{10}\\[3pt]
0&\tfrac1{15}&\tfrac1{15}}=(AB)^{-1}.
\]
\vcheck{Row 1 of $AB$ times column 1 of this matrix:
$10(0.9)+16(-0.5)+(-1)(0)=9-8=1$. Row 1 times column 2:
$\tfrac{10}{6}-\tfrac{16}{10}-\tfrac{1}{15}=0$. It is genuinely the inverse of
$AB$.}

\medskip
By contrast, the \emph{other} order gives
\[
A^{-1}B^{-1}=\frac{1}{-150}\mat{-44&\cdots&\cdots\\ \cdots&\cdots&\cdots\\ \cdots&\cdots&\cdots}
\quad\text{with }(1,1)\text{ entry }\frac{-44}{-150}=\frac{22}{75}\approx 0.293,
\]
which is not the $0.9$ required in position $(1,1)$ of $(AB)^{-1}$.

\medskip
\textbf{(b)} \textbf{No.} $A^{-1}B^{-1}\neq (AB)^{-1}$, as the entry above
shows. The two agree only when $A$ and $B$ commute.

\medskip
\textbf{(c)} \textbf{Yes.} $(AB)^{-1}=B^{-1}A^{-1}$ always, whenever both
inverses exist. The one-line proof:
\[
(AB)(B^{-1}A^{-1})=A(BB^{-1})A^{-1}=AIA^{-1}=AA^{-1}=I .
\]
The order reverses for the same reason it does for transposes.

\medskip
\textbf{(d) Does $A+B$ have an inverse?}
\[
A+B=\mat{0&1&-1\\8&7&4\\-4&-2&4},\qquad
\det(A+B)=0-1(32+16)-1(-16+28)=-48-12=-60\neq 0,
\]
so yes. With
$\operatorname{adj}(A+B)=\mat{36&-2&11\\-48&-4&-8\\12&-4&-8}$,
\[
(A+B)^{-1}=\frac{1}{-60}\mat{36&-2&11\\-48&-4&-8\\12&-4&-8}
=\mat{-\tfrac35&\tfrac1{30}&-\tfrac{11}{60}\\[3pt]
\tfrac45&\tfrac1{15}&\tfrac2{15}\\[3pt]
-\tfrac15&\tfrac1{15}&\tfrac2{15}}.
\]
\vcheck{Row 2 of $A+B$ times column 2:
$\tfrac{8}{30}+\tfrac{7}{15}+\tfrac{4}{15}=\tfrac{4+7+4}{15}=1$.}

\medskip
\textbf{(e) Does $A-B$ have an inverse?}
\[
A-B=\mat{2&5&-1\\2&-3&2\\0&4&-4},\qquad
\det(A-B)=2(12-8)-5(-8-0)-1(8-0)=8+40-8=40\neq 0,
\]
so yes. With
$\operatorname{adj}(A-B)=\mat{4&16&7\\8&-8&-6\\8&-8&-16}$,
\[
(A-B)^{-1}=\frac{1}{40}\mat{4&16&7\\8&-8&-6\\8&-8&-16}
=\mat{\tfrac1{10}&\tfrac25&\tfrac7{40}\\[3pt]
\tfrac15&-\tfrac15&-\tfrac3{20}\\[3pt]
\tfrac15&-\tfrac15&-\tfrac25}.
\]
\vcheck{Row 3 of $A-B$ times column 3:
$0(\tfrac{7}{40})+4(-\tfrac{3}{20})+(-4)(-\tfrac25)=-0.6+1.6=1$.}

\ans{(a) $A^{-1}=\tfrac{1}{30}\mat{3&1&-11\\6&2&8\\-9&7&13}$,
$B^{-1}=\tfrac15\mat{23&8&-2\\-14&-4&1\\1&1&1}$, and
$(AB)^{-1}=B^{-1}A^{-1}$ as displayed. \quad
(b) No. \quad (c) Yes, always. \quad
(d) Yes, $\det(A+B)=-60$. \quad (e) Yes, $\det(A-B)=40$.}
\end{solution}

\begin{pitfall}
The adjoint is the \textbf{transpose} of the cofactor matrix. Forgetting the
transpose produces a matrix that fails $AA^{-1}=I$ in exactly the off-diagonal
entries --- which is why the check above tests an off-diagonal entry, not just
a diagonal one.
\end{pitfall}

%---------------------------------------------------------------------
\section{Four inverses by cofactors}

\begin{question}{Q.3}
Find the inverse, if it exists, of the given matrices using both the cofactor
approach and Gaussian reduction.
\[
\text{(a) }\mat{6&0&-1\\1&3&5\\-3&1&0}\quad
\text{(b) }\mat{1&2&3\\1&3&5\\1&5&12}\quad
\text{(c) }\mat{-1&2&-3\\0&1&5\\2&4&7}\quad
\text{(d) }\mat{8&0&-2\\1&-2&1\\-3&1&1}
\]
\end{question}

\begin{solution}
The cofactor route is shown in full for (a) and summarised for the rest; the
Gaussian route --- row-reducing $[A\mid I]$ to $[I\mid A^{-1}]$ --- gives
identical answers and is sketched at the end.

\medskip
\textbf{(a)} $\det=6(0-5)-0+(-1)(1+9)=-30-10=-40$.
Cofactors:
\[
\begin{array}{lll}
C_{11}=-5, & C_{12}=-15, & C_{13}=10,\\
C_{21}=-1, & C_{22}=-3,  & C_{23}=-6,\\
C_{31}=3,  & C_{32}=-31, & C_{33}=18 .
\end{array}
\]
Transposing and dividing by $\det$:
\[
A^{-1}=\frac{1}{-40}\mat{-5&-1&3\\-15&-3&-31\\10&-6&18}
=\frac{1}{40}\mat{5&1&-3\\15&3&31\\-10&6&-18}.
\]
\vcheck{Row 1 $(6,0,-1)$ times column 1
$\tfrac{1}{40}(5,15,-10)^{t}$: $\tfrac{30+0+10}{40}=1$. Row 2 $(1,3,5)$ times
the same column: $\tfrac{5+45-50}{40}=0$.}

\medskip
\textbf{(b)} $\det=1(36-25)-2(12-5)+3(5-3)=11-14+6=3$.
\[
A^{-1}=\frac13\mat{11&-9&1\\-7&9&-2\\2&-3&1}.
\]
\vcheck{Row 3 $(1,5,12)$ times column 1 $\tfrac13(11,-7,2)^{t}$:
$\tfrac{11-35+24}{3}=0$. Row 1 times column 1: $\tfrac{11-14+6}{3}=1$.}

\medskip
\textbf{(c)} $\det=-1(7-20)-2(0-10)+(-3)(0-2)=13+20+6=39$.
\[
A^{-1}=\frac{1}{39}\mat{-13&-26&13\\10&-1&5\\-2&8&-1}.
\]
\vcheck{Row 1 $(-1,2,-3)$ times column 1 $\tfrac1{39}(-13,10,-2)^{t}$:
$\tfrac{13+20+6}{39}=1$. Row 2 $(0,1,5)$ times the same: $\tfrac{0+10-10}{39}=0$.}

\medskip
\textbf{(d)} $\det=8(-2-1)-0+(-2)(1-6)=-24+10=-14$.
\[
A^{-1}=\frac{1}{-14}\mat{-3&-2&-4\\-4&2&-10\\-5&-8&-16}
=\frac{1}{14}\mat{3&2&4\\4&-2&10\\5&8&16}.
\]
\vcheck{Row 1 $(8,0,-2)$ times column 1 $\tfrac1{14}(3,4,5)^{t}$:
$\tfrac{24+0-10}{14}=1$. Row 3 $(-3,1,1)$ times the same:
$\tfrac{-9+4+5}{14}=0$.}

\medskip
\textbf{The Gaussian route, on (b).} Augment with the identity and reduce:
\[
\left[\begin{array}{ccc|ccc}
1&2&3&1&0&0\\ 1&3&5&0&1&0\\ 1&5&12&0&0&1
\end{array}\right]
\xrightarrow[R_3-R_1]{R_2-R_1}
\left[\begin{array}{ccc|ccc}
1&2&3&1&0&0\\ 0&1&2&-1&1&0\\ 0&3&9&-1&0&1
\end{array}\right]
\]
\[
\xrightarrow{R_3-3R_2}
\left[\begin{array}{ccc|ccc}
1&2&3&1&0&0\\ 0&1&2&-1&1&0\\ 0&0&3&2&-3&1
\end{array}\right]
\xrightarrow{\tfrac13 R_3}
\left[\begin{array}{ccc|ccc}
1&2&3&1&0&0\\ 0&1&2&-1&1&0\\ 0&0&1&\tfrac23&-1&\tfrac13
\end{array}\right],
\]
and back-substituting ($R_2-2R_3$, then $R_1-3R_3$, then $R_1-2R_2$) yields
\[
\left[\begin{array}{ccc|ccc}
1&0&0&\tfrac{11}{3}&-3&\tfrac13\\
0&1&0&-\tfrac73&3&-\tfrac23\\
0&0&1&\tfrac23&-1&\tfrac13
\end{array}\right],
\]
which is $\tfrac13\mat{11&-9&1\\-7&9&-2\\2&-3&1}$ --- the cofactor answer.

\ans{(a) $\tfrac{1}{40}\mat{5&1&-3\\15&3&31\\-10&6&-18}$ \quad
(b) $\tfrac13\mat{11&-9&1\\-7&9&-2\\2&-3&1}$ \quad
(c) $\tfrac{1}{39}\mat{-13&-26&13\\10&-1&5\\-2&8&-1}$ \quad
(d) $\tfrac{1}{14}\mat{3&2&4\\4&-2&10\\5&8&16}$. All four determinants are
non-zero, so all four inverses exist.}
\end{solution}

%---------------------------------------------------------------------
\section{Systems by matrices and by Cramer's rule}

\begin{question}{Q.4}
Solve the following systems of linear equations using matrices as well as
Cramer's rule.
\[
\text{(a) }
\begin{aligned}
2x-y+z&=5\\ 4x+2y+3z&=8\\ 3x-4y-z&=3
\end{aligned}
\qquad
\text{(b) }
\begin{aligned}
x+y\ \ \ &=2\\ 2x\ \ \ -z&=1\\ 2y-3z&=-1
\end{aligned}
\qquad
\text{(c) }
\begin{aligned}
2x+y+3z&=3\\ x+y-2z&=0\\ -3x-y+2z&=-4
\end{aligned}
\]
\end{question}

\begin{solution}
\textbf{(a)} The coefficient determinant, expanded along row 1:
\[
D=\dt{2&-1&1\\4&2&3\\3&-4&-1}
=2(-2+12)+1(-4-9)+1(-16-6)=20-13-22=-15 .
\]
$D\neq 0$, so Cramer's rule applies. Replace one column at a time by the
right-hand side $(5,8,3)^{t}$:
\[
D_x=\dt{5&-1&1\\8&2&3\\3&-4&-1}=5(10)+1(-17)+1(-38)=-5,
\]
\[
D_y=\dt{2&5&1\\4&8&3\\3&3&-1}=2(-17)-5(-13)+1(-12)=19,
\]
\[
D_z=\dt{2&-1&5\\4&2&8\\3&-4&3}=2(38)+1(-12)+5(-22)=-46 .
\]
\[
x=\frac{D_x}{D}=\frac{-5}{-15}=\frac13,\qquad
y=\frac{D_y}{D}=\frac{19}{-15}=-\frac{19}{15},\qquad
z=\frac{D_z}{D}=\frac{-46}{-15}=\frac{46}{15}.
\]
\vcheck{Equation 1: $2(\tfrac13)+\tfrac{19}{15}+\tfrac{46}{15}
=\tfrac{10+19+46}{15}=\tfrac{75}{15}=5$. \
Equation 3: $1+\tfrac{76}{15}-\tfrac{46}{15}=1+2=3$. Both hold.}

\medskip
\textbf{(b)} Write the missing variables in explicitly --- the zeros are what
make this system easy:
\[
A=\mat{1&1&0\\2&0&-1\\0&2&-3},\qquad b=\mat{2\\1\\-1}.
\]
\[
D=1(0+2)-1(-6-0)+0=2+6=8 .
\]
\[
D_x=\dt{2&1&0\\1&0&-1\\-1&2&-3}=8,\qquad
D_y=\dt{1&2&0\\2&1&-1\\0&-1&-3}=8,\qquad
D_z=\dt{1&1&2\\2&0&1\\0&2&-1}=8 .
\]
\[
x=\frac88=1,\qquad y=\frac88=1,\qquad z=\frac88=1 .
\]
\vcheck{$1+1=2$; \ $2(1)-1=1$; \ $2(1)-3(1)=-1$. All three hold.}

\medskip
\textbf{(c)}
\[
D=\dt{2&1&3\\1&1&-2\\-3&-1&2}=2(2-2)-1(2-6)+3(-1+3)=0+4+6=10 .
\]
\[
D_x=\dt{3&1&3\\0&1&-2\\-4&-1&2}=20,\qquad
D_y=\dt{2&3&3\\1&0&-2\\-3&-4&2}=-16,\qquad
D_z=\dt{2&1&3\\1&1&0\\-3&-1&-4}=2 .
\]
\[
x=\frac{20}{10}=2,\qquad y=\frac{-16}{10}=-\frac85,\qquad
z=\frac{2}{10}=\frac15 .
\]
\vcheck{Equation 1: $4-\tfrac85+\tfrac35=4-1=3$. \
Equation 2: $2-\tfrac85-\tfrac25=2-2=0$. \
Equation 3: $-6+\tfrac85+\tfrac25=-6+2=-4$. All three hold.}

\medskip
\textbf{The matrix method gives the same thing.} $Ax=b \Rightarrow x=A^{-1}b$.
For (c), $\det A=10$ and
$A^{-1}=\tfrac{1}{10}\operatorname{adj}A$; multiplying by
$b=(3,0,-4)^{t}$ reproduces $(2,\,-\tfrac85,\,\tfrac15)^{t}$. Cramer's rule is
simply this computation with the common denominator $\det A$ kept visible,
which is why it is faster when only one variable is wanted.

\ans{(a) $x=\tfrac13$, $y=-\tfrac{19}{15}$, $z=\tfrac{46}{15}$ \quad
(b) $x=y=z=1$ \quad
(c) $x=2$, $y=-\tfrac85$, $z=\tfrac15$}
\end{solution}

\begin{pitfall}
In system (b), the equations are written with variables missing. Insert the
zero coefficients ($x+y+0z=2$, $2x+0y-z=1$, $0x+2y-3z=-1$) before forming
$D$. Reading the coefficients straight off the printed equations puts numbers
in the wrong columns and is the commonest way to lose all the marks on an
otherwise easy question.
\end{pitfall}

%---------------------------------------------------------------------
\section{An electronics production plan --- and a defective question}
\label{sec:u6q5}

\begin{question}{Q.5}
An electronic company manufactures resistors, transistors and computer chips.
Each resistor requires 3 units of copper, 2 units of zinc and 1 unit of glass.
Each transistor requires 3, 2 and 1 units of each material respectively, and
each computer chip requires 2, 1 and 2 units of these materials respectively.
Using 800 units of copper, 400 units of zinc and 500 units of glass, how many
of each item can be prepared?
\end{question}

\begin{solution}
Let $r$, $t$ and $c$ be the numbers of resistors, transistors and chips. Reading
the requirements off the question:
\[
\begin{aligned}
\text{copper}&: & 3r+3t+2c&=800\\
\text{zinc}  &: & 2r+2t+ c&=400\\
\text{glass} &: & \ r+\ t+2c&=500
\end{aligned}
\qquad
A=\mat{3&3&2\\2&2&1\\1&1&2}.
\]

\textbf{The coefficient matrix is singular.} Columns 1 and 2 of $A$ are
identical, because the question gives resistors and transistors the
\emph{same} requirement $(3,2,1)$. Therefore
\[
\det A=0,
\]
and Cramer's rule cannot be applied at all. Structurally, the two products are
mathematically indistinguishable: no amount of material data can tell a
resistor from a transistor, so $r$ and $t$ can never be determined separately
--- only their sum.

\medskip
\textbf{And the system is inconsistent.} Put $s=r+t$. The three equations become
\[
3s+2c=800,\qquad 2s+c=400,\qquad s+2c=500 .
\]
From the second, $c=400-2s$. Substituting into the first,
\[
3s+2(400-2s)=800 \;\Longrightarrow\; 3s+800-4s=800
\;\Longrightarrow\; -s=0 \;\Longrightarrow\; s=0,
\]
so $c=400$. Testing these in the third equation:
\[
s+2c=0+800=800 \neq 500 .
\]
The third equation is violated, so no $(r,t,c)$ satisfies all three.

\ans{As printed, the question has \textbf{no solution}. The coefficient matrix
is singular ($\det A=0$, columns 1 and 2 identical, because resistors and
transistors are given identical material requirements), and the reduced system
in $s=r+t$ and $c$ is inconsistent: it forces $s=0,\ c=400$, which fails the
glass constraint $500$.}

\medskip
\textbf{The repaired problem.} The exercise is clearly intended to have
distinct requirements, and the smallest edit that achieves it is to change the
transistor's zinc requirement from 2 to 1 --- so a transistor needs $3$ copper,
$1$ zinc, $1$ glass. The system becomes
\[
\begin{aligned}
3r+3t+2c&=800\\
2r+\ t+\ c&=400\\
\ r+\ t+2c&=500
\end{aligned}
\qquad
A=\mat{3&3&2\\2&1&1\\1&1&2},
\]
\[
D=3(2-1)-3(4-1)+2(2-1)=3-9+2=-4 \neq 0,
\]
so Cramer's rule now runs normally:
\[
D_r=\dt{800&3&2\\400&1&1\\500&1&2}=-300,\quad
D_t=\dt{3&800&2\\2&400&1\\1&500&2}=-300,\quad
D_c=\dt{3&3&800\\2&1&400\\1&1&500}=-700,
\]
\[
r=\frac{-300}{-4}=75,\qquad t=\frac{-300}{-4}=75,\qquad c=\frac{-700}{-4}=175 .
\]
\vcheck{Copper: $3(75)+3(75)+2(175)=225+225+350=800$. \
Zinc: $2(75)+75+175=400$. \
Glass: $75+75+2(175)=500$. All three materials are consumed exactly.}

\medskip
The \emph{method} --- form $A$, test $\det A$, apply Cramer --- is the
examinable content, and it is unaffected. The printed data simply is not
solvable.
\end{solution}

\begin{pitfall}
Compute $D$ \emph{first}, always. If $D=0$, Cramer's rule is unavailable and
the correct answer is a statement about the system --- ``no solution'' or
``infinitely many'' --- not a set of numbers. Producing numbers from a singular
system means an arithmetic error has been made somewhere.
\end{pitfall}

%---------------------------------------------------------------------
\section{Machine time and material cost}

\begin{question}{Q.6}
In a pottery factory, a machine takes 3 minutes to make a glass and 2 minutes to
make a plate. The material for a glass costs \$0.25 and for a plate \$0.20. If
the machine runs for 10 hours and exactly \$50 is used on material, how many
glasses and plates can be produced?
\end{question}

\begin{solution}
Let $g$ be the number of glasses and $p$ the number of plates. Ten hours is
$600$ minutes, so:
\[
\begin{aligned}
\text{time}    &: & 3g+2p&=600\\
\text{material}&: & 0.25g+0.20p&=50
\end{aligned}
\]
Clear the decimals in the second equation by multiplying by 20:
\[
5g+4p=1000 .
\]

\textbf{By elimination.} Double the time equation so the $p$ terms match:
\[
6g+4p=1200 .
\]
Subtracting the material equation:
\[
(6g+4p)-(5g+4p)=1200-1000 \;\Longrightarrow\; g=200 .
\]
Back into $3g+2p=600$:
\[
600+2p=600 \;\Longrightarrow\; p=0 .
\]

\vcheck{Time: $3(200)+2(0)=600$ minutes $=10$ hours. \
Material: $0.25(200)+0.20(0)=\$50$. Both constraints are met exactly.}

\medskip
\textbf{Reading the answer.} The solution sits on the boundary of the feasible
region: the two constraints are satisfied only by making glasses alone. The
system is not degenerate --- $D=\dt{3&2\\5&4}=12-10=2\neq0$, so the solution is
unique --- it simply happens that the two constraints bind simultaneously at a
corner where $p=0$.

\ans{$g=200$ glasses and $p=0$ plates. Running for 10 hours on exactly \$50 of
material is only possible by producing 200 glasses and no plates.}
\end{solution}
